\section{Factorisation in Polynomial Rings}

Recall that is \(K\) is a field, then \(K\brac{x}\) is Euclidean (recall \autoref{prop:euclidean-algorithm-polynomials}). Therefore, we have
\begin{itemize}
    \item any ideal \(I \idealsub K\brac{x}\) is principle (so \(I = \para{f}\)), by \autoref{prop:euclidean-pid};
    \item \(f \in K\brac{x}\) is irreducible if and only if it is prime, by \autoref{lem:irreducible-primes-pid} and \autoref{prop:primes-are-irreducible};
    \item let \(f\) be irreducible, and we consider \(\para{f} \idealsub K\brac{x}\). Suppose \(\para{f} \subseteq J \idealsub k\brac{x}\), then \(J = \para{g}\), so \(f = gh\). But either \(g\) and \(f\) are associates, or \(g\) is a unit. So \(\para{f}\) must be maximal. In fact, this is true for all PIDs.
\end{itemize}

In particular, for \(f \in K\brac{x}\),
\[
    f \text{ is irreducible} \iff K\brac{x} / \para{f} \text{ is a field}.
\]

\begin{definition}[Content of a Polynomial, Primitive Polynomials]
    \label{def:content-polynomial-primitive}%
    Let \(R\) be a UFD, and \(f = a_0 + a_1 x + a_2 x^2 + \cdots + a_n x^n\in R\brac{x}\). The \emph{content} of \(f\) is
    \[
        c\para{f} = \gcd\brce{a_0, \cdots, a_n} \in R.
    \]

    If \(c\para{f}\) is a unit, then we say \(f\) is \emph{primitive}.
\end{definition}

\begin{theorem}[Gauss's Lemma]
    \label{thm:gauss-lemma}%
    Let \(R\) be a UFD, and \(f \in R\brac{x}\). If \(f\) is primitive, then \(f\) is irreducible in \(R \brac{x}\) if and only if \(f\) is reducible in \(F\brac{x}\), where \(F\) is a fraction field over \(R\).
\end{theorem}

Recall \autoref{def:field-of-fractions} for fraction fields.

\begin{remark}
    For example, in \(\para{2x + 2} \in \ZZ\brac{x}\) is \emph{not} primitive, since \(\gcd\brce{2, 2} = 2\) is not a unit, and is reducible. However, \(\para{2x + 2} \in \QQ\brac{x}\) \emph{is} primitive, but is irreducible.
\end{remark}

\begin{example}[Gauss's Lemma in \(\ZZ\brac{x}\) and \(\QQ\brac{x}\)]
    \label{ex:application-gausss-lemma}%
    We consider the polynomial
    \[
        x^3 + x + 1 \in \ZZ\brac{x}.
    \]

    The content of this polynomial is \(1\). We show this is irreducible in \(\ZZ\brac{x}\) and hence in \(\QQ\brac{x}\).

    If it were reducible, we let
    \[
        x^3 + x + 1 = gh
    \]
    with \(g, h \in \ZZ\brac{x}\). Therefore, they must have degrees of at least \(1\) (using the content of the polynomials). Therefore, we choose \(g = b_0 + b_1 x\) and \(h = c_0 + c_1 x + c_2 x^2\).

    Since \(x^2 + x + 1 = gh\), we have \(b_0 c_0 = 1\), \(b_1 c_2 = 1\), and so \(b_0\) and \(b_1\) are \(\pm 1\).
    
    This means \(\pm 1\) is a root of \(x^3 + x + 1\). But it is not by putting it into the polynomial!

    Therefore, \(x^3 + x + 1\) is irreducible in \(\ZZ\brac{x}\), and hence in \(\QQ\brac{x}\) in \autoref{thm:gauss-lemma}.
\end{example}

\begin{lemma}[Gauss's Lemma on Primitivity]
    \label{lem:gauss-lemma-primitive}
    Let \(R\) be a UFD, and \(f, g \in R\brac{x}\). If \(f\) and \(g\) are primitive, then \(fg\) is primitive.
\end{lemma}

\begin{proof}
    Let \(f = a_0 + a_1 x + \cdots + a_n x^n\) and \(g = b_0 + b_1 x + \cdots + b_m x^m\), with \(a_n, b_m \neq 0\) and both primitive.

    Now, suppose \(c\para{fg}\) is not a unit, we may find some \(p \in R\), irreducible, with \(p \divides c\para{fg}\). By the setup, we have \(p\notdivides c\para{f}\), and \(p\notdivides c\para{g}\).

    Suppose \(k\) is the smallest integer such that \(p \notdivides a_k\) and \(l\) is the smallest such that \(p \notdivides b_l\). We consider the coefficient of \(x^{k + l}\), which is given by
    \[
        \sum_{i + j = k + l} a_i b_j.
    \]

    By assumption, \(p\) divides this coefficient, and \(p\) divides the elements
    \[
        a_{k - 1} b_{l + 1}, a_{k - 2} b_{l + 2}, \ldots, a_0 b_{l + k},
    \]
    and
    \[
        a_{k + 1} b_{l - 1}, a_{k + 2} b_{l - 2}, \ldots, a_{l + k} b_0
    \]
    by the assumption of the minimality of \(l\) and \(k\). Therefore, \(p \divides a_k b_l\) which gives a contradiction. 
\end{proof}

\begin{corollary}[Contents are Multiplicative up to Units]
    \label{cor:content-multiplicative}%
    Let \(R\) be a UFD and \(f, g \in R\brac{x}\). Then \(c\para{fg}\) and \(c\para{f} \cdot c\para{g}\) are associates.
\end{corollary}

\begin{proof}
    Let
    \[
        f = c\para{f} f_1, \quad g = c\para{g} g_1
    \]
    for \(f_1, g_1 \in R\brac{x}\), primitive, then \(fg = c\para{f} c\para{g} f_1 g_1\), but \(f_1 g_1\) is primitive by \autoref{lem:gauss-lemma-primitive}.

    Since \(f_1, g_1\) are primitive, \(c\para{f} \cdot c\para{g}\) is indeed a gcd for the coefficients of \(fg\).
\end{proof}

\begin{proof}[of \autoref{thm:gauss-lemma} (Gauss's Lemma)]
    If \(f = gh\) in \(R\brac{x}\) with \(g, h\) non-units, since \(f\) is primitive, we have \(g, h\) are also primitive. Therefore, \(\deg \para{g}, \deg\para{h} > 0\). Therefore, \(f\) is reducible in \(F\brac{x}\).

    Conversely, let \(f = gh\) in \(F\brac{x}\) with \(g, h\) being non-units, so \(\deg\para{g}, \deg\para{h} > 0\). We may find \(a, b\) such that \(\para{ab}f = \para{ag} \para{bh}\), where \(\para{ag}, \para{bh} \in R\brac{x}\). (This is guaranteed by the definition, see \autoref{def:field-of-fractions}.)

    Now, we can write \(\para{ag} = c\para{ag} g_1\) and \(\para{bh} = c\para{bh} h_1\), with \(g, h\) primitive.

    We know since \(f\) is primitive, we have \(c\para{abf} = ab\). By \autoref{cor:content-multiplicative}, we have
    \[
        ab = c\para{abf} = c\para{\para{ag} \cdot \para{bh}} = c\para{ag} \cdot c\para{bh} \cdot u
    \]
    where \(u\) is a unit. Also,
    \[
        abf = c\para{ag} \cdot c\para{bh} g_1 h_1 = u\inverse ab g_1 h_1,
    \]
    and we cancel \(ab\), giving \(f = u\inverse g_1 h_1\), finishing the proof.
\end{proof}

\begin{proposition}[Divisibility in Polynomial Ring over Field of Fractions]
    \label{prop:divisibility-polynomial-ring-fof}%
    Let \(R\) be a UFD, and \(F\) be a field of fractions of \(R\). Let \(g \in R\brac{x}\) be primitive. If \(f\) is divisible by \(g\) in \(F\brac{x}\), then it is also true in \(R\brac{x}\).

    In other words, let
    \[
        I = \para{g} \idealsub R\brac{x}, \quad J = \para{g} \idealsub F\brac{x},
    \]
    then \(I = J \cap R\brac{x}\).
\end{proposition}

\begin{proof}
    Let \(f \in J \cap R\brac{x}\). Then we can write \(f = gh\) for some \(h \in F\brac{x}\). Multiplying by some \(b \in R\), we get
    \[
        bf = g \para{bh}
    \]
    now for \(\para{bh} \in R\brac{x}\). We now see
    \[
        \para{bh} = c\para{bh} \cdot h_1
    \]
    with \(h_1\) primitive. Since \(g\) is primitive, we see \(g h_1\) is primitive, so
    \[
        c\para{bh} = u \cdot c\para{bf}
    \] 
    where \(u\) is a unit. But \(c\para{bf} = b \cdot c\para{f}\), since \(b, f \in R\brac{x}\). Now, observe
    \[
        bf = g u b \cdot c\para{f} \cdot h_1
    \]
    and cancelling out the \(b\) gives divisibility as desired.
\end{proof}

\begin{theorem}[Polynomial Ring over UFD are UFDs]
    \label{thm:poynomial-over-ufd-ufd}%
    If \(R\) is a UFD, then \(R\brac{x}\) is a UFD.
\end{theorem}

\begin{proof}
    Let \(f \in R\brac{x}\). We let \(f = c\para{f} \cdot f_1\).

    We may first factorise the content as
    \[
        c\para{f} = p_1 \cdot p_2 \cdots p_n
    \]
    uniquely, since \(R\) is a UFD. This reduces now to proving a factorisation exists in \(f_1\), where \(f_1\) is primitive.

    If \(f_1\) is irreducible, then we are done. If not, then \(f_1 = f_{11} \cdot f_{12}\) where \(f_{11}\) and \(f_{12}\) are of smaller degrees, and primitive. We continue doing the same with \(f_{11}\) and \(f_{12}\), and the process terminates until the degree drops at some step.
    
    Therefore, we may factorise
    \[
        f_1 = q_1 \cdot q_2 \cdots q_m
    \]
    into irreducibles, and hence \(f\) as
    \[
        f_1 = p_1 \cdots p_n \cdot q_1 \cdots q_m
    \]

    It remains to show the uniqueness of this factorisation. The uniqueness for the \(p_i\)s follows immediately from \(R\) being a UFD, by cancelling content, up to units in \(R\). The uniqueness of the \(q_i\) is as follows. Consider two factorisations
    \[
        f_1 = q_1 \cdots q_m = r_1 \cdots r_l
    \]
    in \(R\brac{x}\).

    We can view this in \(F\brac{x}\). But \(F\brac{x}\) is Euclidean, and hence a PID (by \autoref{prop:euclidean-pid}), and hence a UFD (by \autoref{thm:pid-ufd}). Therefore, \(l = m\), and up to reordering, \(r_i\)s and \(q_i\)s are associates in \(F\brac{x}\).

    But by \autoref{prop:divisibility-polynomial-ring-fof}, they must also be associates in \(R\brac{x}\). This finishes our proof.
\end{proof}

Examples of UFDs now include \(\ZZ\brac{x}, \RR\brac{x, y}, F\brac{x, y}\) for a field \(F\), and \(\ZZ\brac{x_1, \ldots, x_n}\). We may start from any UFD and construct a new UFD by taking the polynomial ring. (But recall that \(\ZZ\brac{\sqrt{-5}}\) is not a UFD.)